%%%%%%%%%%%%
%
% $Autor: Wings $
% $Datum: 2019-03-05 08:03:15Z $
% $Pfad: SetUp.tex $
% $Version: 4250 $
% !TeX spellcheck = en_GB/de_DE
% !TeX encoding = utf8
% !TeX root = manual 
% !TeX TXS-program:bibliography = txs:///biber
%
%%%%%%%%%%%%

\chapter{System Operation}

This chapter outlines how to start, run, and safely shut down the Vehicle License Plate Detection System.

\section{Starting the System}
\begin{enumerate}
	\item \textbf{Power On:} Connect the power adapter to the Jetson Nano and allow it to boot.
	\item \textbf{Navigate to Project Directory:} Open a terminal and navigate to the project folder.
	\item \textbf{Activate Environment:} If using a virtual environment, activate it (e.g., \texttt{source venv/bin/activate}).
\end{enumerate}

\section{Running the Detection Script}
\begin{enumerate}
	\item \textbf{Start Script:} Run the main detection script (e.g., \texttt{python3 detect\_plates.py}).
	\item \textbf{View Output:} The system will process the camera feed, displaying detected license plates and activating the LED indicator if configured.
\end{enumerate}

\section{Shutting Down the System}
\begin{enumerate}
	\item \textbf{Stop Script:} Use \texttt{Ctrl + C} to terminate the script.
	\item \textbf{Power Off:} Shut down the Jetson Nano with \texttt{sudo shutdown now}.
\end{enumerate}

\section{Safety Notes}
Ensure stable setup, avoid extreme conditions, and check cable connections periodically.

This completes the basic operational steps. For troubleshooting guidance, refer to the next chapter.